\section{Scope, exact targets, and conventions}
\label{sec:scope}

The starting point is the manuscript directory and the incoming archive
directory of the ProveIt repository, inspected at the commit printed on
the title page \cite{Repo}. A repository-relative claim of progress is
different from a claim of priority in the mathematical literature.
This article makes the former precise. Classical Gamma quotients,
harmonic Bell polynomials, height-one coefficient series, and
low-weight multiple-zeta reductions are explicitly attributed. The
proofs developed here organize these tools into a fifth-power bridge
and a finite moment reconstruction theorem. The bounded literature
review is not a proof that no equivalent formulation exists elsewhere.

Two especially relevant incoming reports are \emph{Quartic Tails and
Complex Resolvents} and \emph{Barnes Resolvents and Higher Harmonic Identities}
\cite{Quartic,Barnes}. They already establish the quartic digamma bridge,
the all-depth reciprocal-Gamma singular Laurent coefficients, and finite
formulas for two-even-tail moments. Their stated further questions,
rather than results already proved in their bodies, determine our targets.

\begin{center}
\small
\begin{tabular}{@{}>{\raggedright\arraybackslash}p{.22\textwidth}>{\raggedright\arraybackslash}p{.69\textwidth}@{}}
\toprule
Source target & Result and precise boundary \\
\midrule
Quartic report, Q1: fifth digamma power &
Theorem~\ref{q5:main} gives the complete fifth-power finite part, its
mixed coordinate, and all endpoint constants. \\[5pt]
Quartic report, Q2: canonical polynomial &
The scalar residue normalization is proved for every degree, uniquely
up to an additive constant. An exact differential normalization is also
proved. The stronger claim that scalar polynomials remove every
lower-depth arithmetic coordinate remains open. \\[5pt]
Quartic report, Q8: initial-moment generation &
Theorem~\ref{tr:main} proves an explicit finite-prefix bound and rational
reconstruction. For equal orders the first $q$ moments suffice. The
same first-$q$ assertion for all unequal orders remains open. \\[5pt]
Manuscript $S_6$, revised $S_8$ &
No proof is asserted here. Existing numerical residual certificates
retain their original evidentiary status. \\
\bottomrule
\end{tabular}
\end{center}

\subsection{Harmonic and spectral notation}

Put $H_n^{(r)}=\sum_{k=1}^n k^{-r}$ and $H_n=H_n^{(1)}$, with every
empty sum zero. For strict indices, use
\begin{align}
 H_n(s_1,\ldots,s_d)&=
 \sum_{n\ge k_1>\cdots>k_d\ge1}k_1^{-s_1}\cdots k_d^{-s_d},\\
 Z_d(s_1,\ldots,s_d)&=
 \sum_{k_1>\cdots>k_d\ge1}k_1^{-s_1}\cdots k_d^{-s_d}
 \label{eq:strictzeta}
\end{align}
in their convergence regions, followed by analytic continuation where
specified. The elementary harmonic polynomials are
\begin{equation}\label{eq:edef}
 e_j(n)=H_n(\{1\}^j),\qquad e_0(n)=1,
 \qquad \sum_{j\ge0}e_j(n)t^j=\prod_{k=1}^n\left(1+\frac tk\right).
\end{equation}
They must be distinguished from the complete, weakly ordered harmonic
sums, whose generator has factors $(1-t/k)^{-1}$.

The ordinary Stieltjes convention is
\begin{equation}\label{eq:stieltjes}
 \zeta(1+u)=\frac1u+
 \sum_{j\ge0}\frac{(-1)^j\gamma_j}{j!}u^j,
 \qquad \gamma=\gamma_0.
\end{equation}
At depth $d$, define literal regular Laurent coefficients by
\begin{equation}\label{eq:lambdadef}
 \lambda_{d,j}=[u^j]Z_d(1+u,\{1\}^{d-1}),\qquad
 \lambda_d=\lambda_{d,1}.
\end{equation}
The coefficient brackets impose neither a factorial nor an alternating
sign. In particular
\[
 \lambda_1=-\gamma_1,\qquad
 \lambda_2=\eta,\quad\lambda_3=\theta,\quad\lambda_4=\xi.
\]
The ordered germ has a pole of order $d$, not a simple pole for $d>1$.
It is therefore unsafe to identify it with a Hurwitz-zeta germ merely
by calling both sets of coefficients generalized Stieltjes constants.

We use the reciprocal-Gamma coefficients
\begin{equation}\label{eq:cdef}
 \frac1{\Gamma(1+t)}=\sum_{j\ge0}c_jt^j,
 \quad c_0=1,\quad c_1=\gamma,\quad
 c_2=\frac{\gamma^2-\zeta(2)}2,\quad
 c_3=\frac{\gamma^3-3\gamma\zeta(2)+2\zeta(3)}6.
\end{equation}
They are obtained recursively by differentiating the standard
log-Gamma expansion \cite{DLMF}:
\begin{equation}\label{eq:crecurrence}
 n c_n=\gamma c_{n-1}
       +\sum_{r=2}^n(-1)^{r+1}\zeta(r)c_{n-r}.
\end{equation}

\subsection{A finite part is a specified constant coefficient}

For a function meromorphic at zero and smooth on $(0,1]$, throughout
this article
\begin{equation}\label{eq:fpdef}
 \FP\int_0^1f(x)\,dx
 =\CT_{\varepsilon\downarrow0}\int_\varepsilon^1f(x)\,dx.
\end{equation}
The expansion uses powers of the coordinate $\varepsilon$ itself and
$\log\varepsilon$. There is no silently changed cutoff, rescaled
coordinate, or substituted spectral regularization. If
$f(x)=\sum_{k=-M}^{\infty}a_kx^k$ near zero, this definition can be
computed without divergent numerical cancellation as
\begin{equation}\label{eq:fpsubtraction}
 \int_0^1\left(f(x)-\sum_{k=-M}^{-1}a_kx^k\right)dx
       +\sum_{k=-M}^{-2}\frac{a_k}{k+1}.
\end{equation}
The contribution from $a_{-1}/x$ is zero. For a meromorphic $F$,
\begin{equation}\label{eq:fpderivative}
 \FP\int_0^1F'(x)\,dx=F(1)-[x^0]F(x).
\end{equation}
This rule follows directly by differentiating the Laurent series; it
will account for every endpoint term in the fifth-power calculation.

All ordinary integrals below are identified as such. All matrix ranks
in the moment reconstruction are formal rational ranks. Neither
finite-part calculations nor rational elimination require unproved
arithmetic independence of zeta or Stieltjes constants.
